\subsection{The boundary-cosine identity in record 86}

\begin{proposition}[Boundary-cosine identity, independent verification]\ours
Let \(K\subset\mathbb R^2\) be bounded and strictly convex with \(C^2\)
boundary. For distinct \(M,N\in\partial K\), let \(c_M\) be the cosine between
the inward unit normal at \(M\) and the chord directed from \(M\) to \(N\);
define \(c_N\) with the chord directed from \(N\) to \(M\). Then
\[
 \int_{\partial K}\int_{\partial K}c_Mc_N\,ds_M\,ds_N=2\pi |K|.
\]
\end{proposition}

\begin{proof}
Parameterize unoriented lines by an angle \(0\leq\theta<\pi\) and a signed
normal coordinate \(p\). Put
\(v=(\cos\theta,\sin\theta)\), \(n=(-\sin\theta,\cos\theta)\).
A transverse line through \(K\) has endpoints
\(M=pn+av\), \(N=pn+bv\), \(a<b\), and chord length \(\ell=b-a\).
Differentiating the endpoint formulas gives
\[
 n\cdot dM=dp+a\,d\theta,\qquad n\cdot dN=dp+b\,d\theta.
\]
Taking the exterior product and then absolute values yields
\[
 c_Mc_N\,ds_M\,ds_N=\ell\,|dp\,d\theta|.
\]
The map from ordered distinct endpoint pairs to unoriented lines is
two-to-one almost everywhere. Hence the boundary integral is
\(2\int_0^\pi\int_{\mathbb R}\ell(p,\theta)\,dp\,d\theta\).
For each fixed \(\theta\), Fubini in orthogonal coordinates gives
\(\int\ell(p,\theta)\,dp=|K|\). This proves the asserted constant.
Tangencies and the diagonal have zero measure and are handled by limits.
\end{proof}

\paragraph{Independent second check by the divergence theorem.}\ours{}
Let \(m_M\) be the inward normal at \(M\), and let \(\nu_N\) be the
outward normal at \(N\). For \(r=x-M\), define
\[
 F_M(x)=\frac{(m_M\cdot r)r}{|r|^2}.
\]
On the boundary, \(F_M(N)\cdot\nu_N=c_Mc_N\), whereas away from \(M\)
\(\operatorname{div}F_M(x)=m_M\cdot(x-M)/|x-M|^2\).
Removing a disk of radius \(\varepsilon\) about \(M\) causes a flux
error tending to zero, since \(F_M\) is bounded. Therefore the double
boundary integral equals
\[
 \int_K\int_{\partial K}
 \frac{\nu_M\cdot(M-x)}{|M-x|^2}\,ds_M\,dx.
\]
Fubini is valid because the singularity \(|x-M|^{-1}\) is locally
integrable in the planar variable. For every interior \(x\), the inner
integral is the flux of \(\nabla_M\log|M-x|\), hence \(2\pi\).
This independently recovers \(2\pi|K|\) without the chord-coordinate
Jacobian or its multiplicity count.

This is a verification of the recorded formula, not a claim of a new
identity. The classical line measure is discussed in
\citet{santalo1989integral}, \S1.2; the 1953 treatise
\citep{santalo1953introduction} establishes the pre-seminar setting,
but an exact Borovikov publication was not verified. \oursComputed{}
For a circle of radius
\(R\), \(c_Mc_N=\sin^2((t-s)/2)\) and the integral is
\(2\pi^2R^2=2\pi|K|\), an independent normalization check.

\subsection{The corrected nutation polynomial in record 83}

\begin{proposition}[Elimination for the recorded top Lagrangian]\ours
For
\[
 L=\frac12\bigl(\dot\theta^2+\dot\varphi^2+\dot\psi^2
       +2\dot\varphi\dot\psi\cos\theta\bigr)-p\cos\theta,
\]
let \(n=\partial L/\partial\dot\varphi\),
\(e=\partial L/\partial\dot\psi\), and let \(h=T+p\cos\theta\).
On \(|u|<1\), \(u=\cos\theta\), the reduced equation is
\[
 \dot u^2=2pu^3-2hu^2+2(ne-p)u+2h-n^2-e^2.
\]
\end{proposition}

\begin{proof}
Direct differentiation gives
\(n=\dot\varphi+u\dot\psi\),
\(e=\dot\psi+u\dot\varphi\), without the two extraneous factors 2
in the notebook. Solving this two-by-two system gives
\[
 \dot\varphi=\frac{n-eu}{1-u^2},\qquad
 \dot\psi=\frac{e-nu}{1-u^2}.
\]
Therefore
\[
 \dot\varphi^2+\dot\psi^2+2u\dot\varphi\dot\psi
   =\frac{n^2+e^2-2neu}{1-u^2},\qquad
 \dot\theta^2=\frac{\dot u^2}{1-u^2}.
\]
Substitution into the energy equation and multiplication by
\(2(1-u^2)\) give
\(\dot u^2=2(h-pu)(1-u^2)-n^2-e^2+2neu\), whose expansion is the claim.
\end{proof}

The physical branch must satisfy \(-1\leq u\leq1\) and a nonnegative
right-hand side. Distinct roots give an elliptic quadrature; repeated roots
give degenerate limits. This checks the specific normalized spherical-top
model written in the notes; it does not silently substitute the general
two-inertia symmetric-top model of the classical literature
\citep{klein1897top}.

\subsection{Testing Gelfand's slice-area formula}

\begin{example}[A fixed-direction counterexample]\oursComputed
For the ellipsoid
\(x^2/a^2+y^2/b^2+z^2/c^2\leq1\), a section at height \(t\) has area
\(S(t)=\pi ab(1-t^2/c^2)\), \(|t|\leq c\). Consequently
\[
 \int_{-c}^{c}|S'(t)|^2\,dt=\frac{8\pi^2a^2b^2}{3c},\qquad
 V=\frac{4\pi abc}{3},\qquad
 \frac{\int|S'|^2}{V}=\frac{2\pi ab}{c^2}.
\]
The ratio depends on shape, so no universal multiplier can make the
single-direction formula valid for all ellipsoids. If the ambiguous source
were read as \(S^2\) instead of \((S')^2\), even homothetic scaling would
rule out a universal constant relating that integral to volume.
\end{example}

\begin{proposition}[An averaged repair, not an attribution]\ours
For the hyperplane Radon transform in \(\mathbb R^3\), ordinary surface
measure on \(S^2\), and Euclidean area on each plane,
\[
 \int_{S^2}\int_{\mathbb R}|\partial_t Rf(\omega,t)|^2\,dt\,d\omega
       =8\pi^2\|f\|_{L^2(\mathbb R^3)}^2.
\]
This holds first for Schwartz functions and extends by density in the
corresponding norm; for indicators of bounded domains of finite volume it
gives an averaged volume formula.
\end{proposition}

\begin{proof}
Use \(\widehat f(\xi)=\int e^{-ix\cdot\xi}f(x)\,dx\).
Fubini gives the Fourier-slice identity
\(\mathcal F_t Rf(\omega,\sigma)=\widehat f(\sigma\omega)\).
One-dimensional Plancherel, the double coverage of polar coordinates by
\((\omega,\sigma)\in S^2\times\mathbb R\), and three-dimensional
Plancherel give, successively,
\[
 \frac1{2\pi}\int_{S^2}\int_{\mathbb R}
       \sigma^2|\widehat f(\sigma\omega)|^2\,d\sigma\,d\omega
 =\frac2{2\pi}\int_{\mathbb R^3}|\widehat f(\xi)|^2\,d\xi
 =8\pi^2\|f\|_2^2.
\]
\end{proof}

Neither the angular integral nor its normalization is written in P000138.
The valid averaged identity therefore cannot by itself establish what
Gelfand intended by \(M\), especially since the following line explicitly
calls \(M\) a universal constant.
